% ============================================
% UPDATE FILE - RAG Workflow & Fixed Figures
% ============================================
% This file contains:
% 1. RAG workflow section (new content)
% 2. Fixed architecture diagram (Figure 1)
% 3. Fixed CV parsing workflow (Figure 2)
% 4. Fixed semantic matching workflow (Figure 3)
% 5. New RAG workflow diagram (Figure 4)

% ============================================
% SECTION: RAG Workflow (ADD TO SECTION IV)
% ============================================
% Insert this after Section IV.C (Semantic Job Matching)

\subsection{Retrieval-Augmented Generation (RAG) for Job Insights}

To provide contextual and accurate job market insights, TalentBridge implements a Retrieval-Augmented Generation (RAG) pipeline. This approach combines vector database retrieval with LLM generation to answer student queries about career paths, skill requirements, and market trends.

\subsubsection{RAG Architecture}

The RAG system consists of five key components as illustrated in Figure~\ref{fig:rag_workflow}:

\begin{enumerate}
\item \textbf{Document Ingestion}: Job descriptions and market data are preprocessed and chunked into semantic units
\item \textbf{Vector Embedding}: Each chunk is embedded using text-embedding-004 (768 dimensions)
\item \textbf{Vector Storage}: Embeddings stored in ChromaDB with metadata (job category, company, salary range)
\item \textbf{Retrieval}: User queries are embedded and matched against the vector database
\item \textbf{Generation}: Retrieved context is injected into Gemini prompts for accurate, grounded responses
\end{enumerate}

\begin{figure}[htbp]
\centering
\begin{tikzpicture}[node distance=1.8cm, scale=0.7, every node/.style={scale=0.7}]

% ===== INDEXING PIPELINE (TOP) =====
\node (docs) [data, minimum width=2.8cm, minimum height=0.9cm] {Job Postings\\(3,237 docs)};
\node (chunk) [process, below of=docs, minimum width=2.8cm] {Chunking\\(3 parts/job)};
\node (embed_doc) [process, below of=chunk, minimum width=2.8cm, fill=accentorange!30] {Embedding\\(768-dim)};
\node (vectordb) [database, below of=embed_doc, yshift=-0.3cm, minimum width=2.8cm] {ChromaDB\\Vector Store};

% ===== QUERY PIPELINE (RIGHT) =====
\node (user_query) [data, right of=docs, xshift=5cm, minimum width=2.8cm, minimum height=0.9cm] {User Query};
\node (embed_query) [process, below of=user_query, minimum width=2.8cm, fill=accentorange!30] {Embed Query\\(768-dim)};
\node (retrieve) [process, below of=embed_query, minimum width=2.8cm, fill=secondarygreen!30] {Similarity\\Search};
\node (top_k) [data, below of=retrieve, minimum width=2.8cm, minimum height=0.9cm] {Top-5\\Contexts};

% ===== GENERATION (BOTTOM) =====
\node (llm) [process, below of=top_k, yshift=-0.5cm, minimum width=2.8cm, fill=primaryblue!30] {Gemini 2.5\\+ Context};
\node (response) [data, below of=llm, minimum width=2.8cm, minimum height=0.9cm] {Grounded\\Response};

% ===== ARROWS =====
% Indexing flow - labels on the right side to avoid overlap
\draw [arrow] (docs) -- node[anchor=west, font=\tiny, xshift=3pt] {raw docs} (chunk);
\draw [arrow] (chunk) -- node[anchor=west, font=\tiny, xshift=3pt] {chunks} (embed_doc);
\draw [arrow] (embed_doc) -- node[anchor=west, font=\tiny, xshift=3pt] {vectors} (vectordb);

% Query flow - no labels to avoid clutter
\draw [arrow] (user_query) -- (embed_query);
\draw [arrow] (embed_query) -- (retrieve);

% Vector DB to Retrieval - label on top to avoid overlap
\draw [arrow] (vectordb) -- node[anchor=south, font=\tiny, yshift=2pt] {} (retrieve);

% Retrieval to Context - label on right side
\draw [arrow] (retrieve) -- node[anchor=west, font=\tiny, xshift=3pt] {top-5} (top_k);

% Context + Query to LLM
\draw [arrow] (top_k) -- (llm);
\draw [arrow] (user_query) -- ++(0,-6.5) -| node[anchor=north, font=\tiny, pos=0.5] {query} (llm);

% LLM to Response
\draw [arrow] (llm) -- (response);

\end{tikzpicture}
\caption{RAG Workflow - Indexing pipeline (left) and query pipeline (right) for grounded job insights}
\label{fig:rag_workflow}
\end{figure}

\subsubsection{Implementation Details}

\textbf{Document Chunking Strategy:}
Each job posting is split into semantic chunks:
\begin{itemize}
\item Job overview and requirements (chunk 1)
\item Detailed responsibilities (chunk 2)
\item Company information and benefits (chunk 3)
\end{itemize}

\textbf{Retrieval Process:}
Given a user query $q$, we:
\begin{enumerate}
\item Generate query embedding: $E_q = \text{Embed}(q)$
\item Retrieve top-k similar chunks: $D = \text{ChromaDB.Query}(E_q, k=5)$
\item Construct augmented prompt: $P = \text{Template}(q, D)$
\item Generate response: $R = \text{Gemini}(P)$
\end{enumerate}

\textbf{Prompt Template:}
\begin{verbatim}
Context: {retrieved_chunks}

User Question: {query}

Based on the job market data above, 
provide a detailed answer. Include 
specific examples and statistics.
\end{verbatim}

\subsubsection{Use Cases}

The RAG system powers three key features:

\begin{enumerate}
\item \textbf{Career Path Guidance}: "What skills do I need for a Data Scientist role?"
\item \textbf{Salary Insights}: "What is the average salary for Frontend Developers in HCMC?"
\item \textbf{Market Trends}: "Which companies are hiring the most in AI/ML?"
\end{enumerate}

\subsubsection{Advantages Over Pure LLM}

\begin{itemize}
\item \textbf{Accuracy}: Grounded in real job data, not hallucinated
\item \textbf{Freshness}: Reflects current market (3,237 recent postings)
\item \textbf{Specificity}: Provides company names, salary ranges, exact requirements
\item \textbf{Transparency}: Can cite source job postings
\end{itemize}

% ============================================
% FIXED FIGURE 1: System Architecture
% ============================================
% Replace the existing Figure 1 with this improved version

\begin{figure*}[t]
\centering
\begin{tikzpicture}[node distance=2cm, scale=0.85, every node/.style={scale=0.85}]

% ===== FRONTEND LAYER =====
\node (frontend_label) [draw=none, fill=none, font=\bfseries] at (0, 0) {Frontend Layer};
\node (frontend_box) [draw, thick, rounded corners, minimum width=14cm, minimum height=2.5cm, fill=primaryblue!10] at (0, -1.5) {};

\node (index) [process, minimum width=2.8cm, minimum height=0.8cm] at (-5, -1.5) {Landing Page};
\node (cvpage) [process, minimum width=2.8cm, minimum height=0.8cm] at (-2, -1.5) {CV Analysis};
\node (jobpage) [process, minimum width=2.8cm, minimum height=0.8cm] at (1.5, -1.5) {Job Listing};
\node (dashboard) [process, minimum width=2.8cm, minimum height=0.8cm] at (5, -1.5) {Dashboard};

% ===== API LAYER =====
\node (api_label) [draw=none, fill=none, font=\bfseries] at (0, -4) {API Layer (FastAPI)};
\node (api_box) [draw, thick, rounded corners, minimum width=14cm, minimum height=2.5cm, fill=secondarygreen!10] at (0, -5.5) {};

\node (upload) [process, minimum width=2.2cm, minimum height=0.7cm, fill=secondarygreen!30] at (-5.5, -5.5) {/upload-cv};
\node (match) [process, minimum width=2.2cm, minimum height=0.7cm, fill=secondarygreen!30] at (-3, -5.5) {/match};
\node (jobs) [process, minimum width=2.2cm, minimum height=0.7cm, fill=secondarygreen!30] at (-0.5, -5.5) {/jobs};
\node (insights) [process, minimum width=2.2cm, minimum height=0.7cm, fill=secondarygreen!30] at (2, -5.5) {/insights};
\node (analytics) [process, minimum width=2.2cm, minimum height=0.7cm, fill=secondarygreen!30] at (4.5, -5.5) {/analytics};

% ===== PROCESSING LAYER =====
\node (processing_label) [draw=none, fill=none, font=\bfseries] at (0, -7.5) {Processing Layer};

% AI Module
\node (ai_box) [draw, thick, rounded corners, minimum width=4cm, minimum height=3cm, fill=accentorange!10] at (-4.5, -9.5) {};
\node (ai_label) [draw=none, fill=none, font=\small\bfseries] at (-4.5, -8.3) {AI Analysis};
\node (gemini) [process, minimum width=3cm, minimum height=0.6cm, fill=accentorange!30, font=\small] at (-4.5, -9.2) {Gemini 2.5 Flash};
\node (embedding) [process, minimum width=3cm, minimum height=0.6cm, fill=accentorange!30, font=\small] at (-4.5, -10.1) {text-embedding-004};

% Vector Module
\node (vector_box) [draw, thick, rounded corners, minimum width=4cm, minimum height=3cm, fill=secondarygreen!10] at (0, -9.5) {};
\node (vector_label) [draw=none, fill=none, font=\small\bfseries] at (0, -8.3) {Vector Search};
\node (chromadb) [database, minimum width=2.5cm, minimum height=1cm, font=\small] at (0, -9.7) {ChromaDB\\768-dim};

% RAG Module (NEW)
\node (rag_box) [draw, thick, rounded corners, minimum width=4cm, minimum height=3cm, fill=primaryblue!10] at (4.5, -9.5) {};
\node (rag_label) [draw=none, fill=none, font=\small\bfseries] at (4.5, -8.3) {RAG Module};
\node (rag_retrieval) [process, minimum width=3cm, minimum height=0.6cm, fill=primaryblue!30, font=\small] at (4.5, -9.2) {Retrieval};
\node (rag_generation) [process, minimum width=3cm, minimum height=0.6cm, fill=primaryblue!30, font=\small] at (4.5, -10.1) {Generation};

% ===== DATA LAYER =====
\node (data_label) [draw=none, fill=none, font=\bfseries] at (0, -12) {Data Layer};
\node (sqlite) [database, minimum width=3cm, minimum height=1.2cm] at (0, -13.2) {SQLite\\6 tables};

% External API
\node (google_api) [process, minimum width=2.5cm, minimum height=1cm, fill=yellow!30] at (-9, -9.5) {Google AI\\API};

% ===== ARROWS =====
% Frontend to API (logical connections)
\draw [arrow] (cvpage) -- (upload);
\draw [arrow] (cvpage) -- (match);
\draw [arrow] (jobpage) -- (jobs);
\draw [arrow] (jobpage) -- (insights);
\draw [arrow] (dashboard) -- (analytics);

% API to Processing (correct routing)
\draw [arrow] (upload) -- (ai_box);
\draw [arrow] (match) -- (vector_box);
\draw [arrow] (insights) -- (rag_box);

% Analytics goes around vector search box to avoid overlap
\draw [arrow] (analytics) -- ++(0,-3) -| ++(2.5,0) |- (sqlite);

% Processing to Data
\draw [arrow] (ai_box) -- (sqlite);
\draw [arrow] (vector_box) -- (chromadb);

% External API connections
\draw [arrow] (google_api) -- (gemini);
\draw [arrow] (google_api) -- (embedding);

% ChromaDB connections
\draw [arrow] (chromadb) -- (rag_retrieval);
\draw [arrow, dashed] (ai_box) -- (chromadb);

\end{tikzpicture}
\caption{TalentBridge System Architecture - Three-tier design with AI-powered backend and RAG module}
\label{fig:architecture}
\end{figure*}

% ============================================
% FIXED FIGURE 2: CV Parsing Workflow
% ============================================
% Replace the existing Figure 2 with this clearer version

\begin{figure}[htbp]
\centering
\begin{tikzpicture}[node distance=1.5cm, scale=0.75, every node/.style={scale=0.75}]

% Main flow
\node (pdf) [data, minimum width=3cm] {PDF File};
\node (extract) [process, below of=pdf, minimum width=3cm] {Extract Text\\(PyPDF2)};
\node (prompt) [process, below of=extract, minimum width=3cm] {Construct\\Prompt};
\node (llm) [process, below of=prompt, minimum width=3cm, fill=accentorange!30] {Gemini 2.5\\Flash};
\node (parse) [process, below of=llm, minimum width=3cm] {Parse JSON\\Response};
\node (validate) [decision, below of=parse, yshift=-0.5cm, minimum width=2.5cm, aspect=2] {Valid?};
\node (store) [database, below of=validate, yshift=-0.8cm, minimum width=2.5cm] {Store to\\SQLite};
\node (index) [database, below of=store, yshift=-0.3cm, minimum width=2.5cm] {Index to\\ChromaDB};

% Retry path - moved further right to avoid overlap
\node (retry) [process, right of=validate, xshift=3.5cm, minimum width=2.8cm] {Retry with\\Different Key};

% Arrows
\draw [arrow] (pdf) -- (extract);
\draw [arrow] (extract) -- node[anchor=west, font=\small] {raw text} (prompt);
\draw [arrow] (prompt) -- (llm);
\draw [arrow] (llm) -- node[anchor=west, font=\small] {JSON} (parse);
\draw [arrow] (parse) -- (validate);
\draw [arrow] (validate) -- node[anchor=west, font=\small] {Yes} (store);
\draw [arrow] (validate) -- node[anchor=north, font=\small, yshift=2pt] {No} (retry);
\draw [arrow] (retry) |- (llm);
\draw [arrow] (store) -- (index);

\end{tikzpicture}
\caption{CV Parsing Workflow with LLM-based extraction and API key rotation}
\label{fig:cv_parsing_flow}
\end{figure}

% ============================================
% FIXED FIGURE 3: Semantic Matching Workflow
% ============================================
% Replace the existing Figure 3 with this improved version

\begin{figure}[htbp]
\centering
\begin{tikzpicture}[node distance=1.5cm, scale=0.75, every node/.style={scale=0.75}]

% Main flow
\node (start) [data, minimum width=3cm] {CV Data};
\node (embed) [process, below of=start, minimum width=3cm, fill=accentorange!30] {Generate\\Embedding};
\node (query) [process, below of=embed, minimum width=3cm, fill=secondarygreen!30] {ChromaDB\\Query};
\node (candidates) [data, below of=query, minimum width=3cm] {Top 50\\Candidates};
\node (filter) [decision, below of=candidates, yshift=-0.5cm, minimum width=2.5cm, aspect=2] {Apply\\Filters?};
\node (rank) [process, below of=filter, yshift=-0.8cm, minimum width=3cm, fill=primaryblue!30] {Gemini AI\\Ranking};
\node (explain) [process, below of=rank, minimum width=3cm] {Generate\\Explanations};
\node (result) [data, below of=explain, minimum width=3cm] {Top-k Jobs\\+ Scores};

% Filter path
\node (apply_filter) [process, right of=filter, xshift=3.5cm, minimum width=2.8cm] {Filter by\\Category/Location};

% Arrows
\draw [arrow] (start) -- (embed);
\draw [arrow] (embed) -- node[anchor=west, font=\small] {768-dim} (query);
\draw [arrow] (query) -- (candidates);
\draw [arrow] (candidates) -- (filter);
\draw [arrow] (filter) -- node[anchor=west, font=\small] {No} (rank);
\draw [arrow] (filter) -- node[anchor=south, font=\small] {Yes} (apply_filter);
\draw [arrow] (apply_filter) |- (rank);
\draw [arrow] (rank) -- (explain);
\draw [arrow] (explain) -- (result);

\end{tikzpicture}
\caption{Semantic Job Matching Workflow with vector similarity and AI ranking}
\label{fig:matching_flow}
\end{figure}

% ============================================
% NOTES FOR INTEGRATION
% ============================================
% 
% TO INTEGRATE INTO MAIN FILE:
% 
% 1. RAG Section:
%    - Insert after Section IV.C (line ~300)
%    - This adds subsection IV.D
% 
% 2. Figure 1 (Architecture):
%    - Replace lines ~93-153 with new version
%    - Now includes RAG module
% 
% 3. Figure 2 (CV Parsing):
%    - Replace lines ~211-251 with new version
%    - Clearer layout, better spacing
% 
% 4. Figure 3 (Matching):
%    - Replace lines ~254-300 with new version
%    - More detailed, shows filter logic
% 
% 5. Figure 4 (RAG):
%    - New figure, insert after Figure 3
%    - Shows complete RAG pipeline
% 
% IMPROVEMENTS:
% - All figures now have consistent styling
% - Better spacing and alignment
% - Clearer labels and annotations
% - Color-coded by function (AI=orange, Vector=green, RAG=blue)
% - Proper node distances for readability
% 
% COMPILE TEST:
% pdflatex update.tex
% (This file is standalone for testing)

